\documentclass[11pt,a4paper]{article}
\usepackage{fontspec,amsmath,amssymb,amsthm,microtype}
\setmainfont{lmroman10-regular.otf}[BoldFont=lmroman10-bold.otf,ItalicFont=lmroman10-italic.otf,BoldItalicFont=lmroman10-bolditalic.otf]
\setmonofont{lmmono10-regular.otf}[Scale=MatchLowercase]
\usepackage[margin=25mm]{geometry}
\usepackage[unicode,hidelinks]{hyperref}
\setlength{\emergencystretch}{3em}
\setlength{\parskip}{0.4em}
\setlength{\parindent}{0pt}
\allowdisplaybreaks[2]
\raggedbottom
\title{\texorpdfstring{Finite-\allowbreak{}dimensional Prékopa–Leindler with Gaussian tails}{Finite-dimensional Prékopa–Leindler with Gaussian tails}}
\author{LangevinLean mathematical exposition}
\date{}
\begin{document}\maketitle
\noindent Source module: \texttt{IPLA/\allowbreak{}Prekopa\allowbreak{}Tensorization.\allowbreak{}lean}. The accompanying HTML pairs this mathematical exposition with the exact Lean statements and records the scope of the formalization.


Positive continuous weights dominated by an integrable Gaussian form a class closed under integrating out one coordinate. This supplies all continuity and integrability facts needed to iterate the one-dimensional transport proof.


\section{\texorpdfstring{Gaussian domination proves integrability and positive integral}{Gaussian domination proves integrability and positive integral}}


Write \(\mathrm{Coordinates}(n)=\mathbb R^n\). A Gaussian weight satisfies continuity, strict positivity, and \(f(x)\le Ce^{-a\sum_i x_i^2}\), with \(C,a>0\). The Euclidean Gaussian integral is finite, and the coordinate identification preserves volume, so domination proves \(f\in L^1\). Positivity on the whole nonzero volume space proves \(\int f>0\). Multiplying by \(c>0\) preserves these properties and changes the envelope constant to \(cC\).




\textbf{Correspondence with the displayed declarations.}


\begin{itemize}

\item \texttt{Gaussian\allowbreak{}Weight}: The class of positive continuous weights with a Gaussian upper bound.
\end{itemize}


{\raggedright\noindent Lean declarations: \texttt{Coordinates}, \texttt{Gaussian\allowbreak{}Weight}, \texttt{integrable}, \texttt{integral\_\allowbreak{}pos}, \texttt{const\_\allowbreak{}mul}.\par}

\section{\texorpdfstring{A coordinate slice and its marginal are integrable and continuous}{A coordinate slice and its marginal are integrable and continuous}}


The insertion map \((s,x)\mapsto(s,x_1,\ldots,x_n)\) is continuous coordinate by coordinate. Define \(f_1(x)=\int_\mathbb R f(s,x)ds\). Since \(\sum_i(s,x)_i^2=s^2+\sum_i x_i^2\ge s^2\), every slice is bounded by \(Ce^{-as^2}\), which is integrable independently of \(x\). Slice continuity supplies measurability and domination proves slice integrability. If \(x_j\to x\), joint continuity gives pointwise convergence in \(s\), and the same Gaussian bound allows dominated convergence. Hence \(f_1\) is continuous.




Write a point in \(\mathbb R^{n+1}\) as \(\operatorname{cons}(z,x)\), where \(z\in\mathbb R\) is the coordinate being integrated. The coordinate-insertion continuity theorem checks its first and remaining coordinates separately. A Gaussian weight has constants \(C,a>0\) satisfying

\[0<f(\operatorname{cons}(z,x))\le C\exp\!\left[-a\left(z^2+\sum_{i=1}^{n}x_i^2\right)\right]\le C e^{-az^2}.\]

The last inequality drops a nonnegative sum of squares, so the majorant is independent of \(x\). It is integrable in \(z\) by the one-dimensional Gaussian integral theorem. The declaration \texttt{slice\_\allowbreak{}integrable} uses positivity to identify the absolute value of the slice with the slice itself, continuity to obtain measurability, and this majorant to prove integrability.



Now let \(x_m\to x\). For every fixed \(z\), continuity of \(f\) and of coordinate insertion gives \(f(\operatorname{cons}(z,x_m))\to f(\operatorname{cons}(z,x))\). The same Gaussian majorant controls every term and the limit. Dominated convergence therefore gives convergence of their real integrals, proving continuity of the first marginal. In the Lean invocation of \texttt{continuous\_\allowbreak{}of\_\allowbreak{}dominated}, the four obligations are slice measurability, the uniform almost-everywhere bound, integrability of that bound, and pointwise continuity in the remaining coordinates. Those are all established before the induction later treats the marginal as another Gaussian weight.




\textbf{Correspondence with the displayed declarations.}


\begin{itemize}

\item \texttt{continuous\_\allowbreak{}cons}: Continuity of inserting the first real coordinate.
\item \texttt{first\allowbreak{}Marginal}: Integrate out the first coordinate of a weight.
\end{itemize}


{\raggedright\noindent Lean declarations: \texttt{continuous\_\allowbreak{}cons}, \texttt{first\allowbreak{}Marginal}, \texttt{slice\_\allowbreak{}integrable}, \texttt{first\allowbreak{}Marginal\_\allowbreak{}continuous}.\par}

\section{\texorpdfstring{Coordinate integration preserves Gaussian tails and total integral}{Coordinate integration preserves Gaussian tails and total integral}}


Retaining the full Gaussian envelope and factoring its exponent gives

\[f_1(x)\le C\left(\int_\mathbb R e^{-as^2}ds\right)e^{-a\sum_i x_i^2}.\]

The new constant is positive and finite, and integrating the positive slice makes \(f_1\) positive everywhere. Together with continuity, this proves that \(f_1\) is again a Gaussian weight. The coordinate splitting equivalence preserves product volume. Its integral formula and Fubini, justified by the Gaussian integrability already proved, then show \(\int f_1=\int f\).



{\raggedright\noindent Lean declarations: \texttt{first\allowbreak{}Marginal\_\allowbreak{}gaussian}, \texttt{integral\_\allowbreak{}first\allowbreak{}Marginal}.\par}

\section{\texorpdfstring{Induct on the dimension}{Induct on the dimension}}


For \(n=0\), the coordinate space is a singleton with volume one; each integral equals its function value and the result is the pointwise premise. For dimension \(n+1\), fix remaining-coordinate vectors \(x,y\). The one-dimensional theorem applied to the slices \(s\mapsto f(s,x)\), \(r\mapsto g(r,y)\), and \(u\mapsto h(u,(1-t)x+ty)\) gives the pointwise Prékopa inequality for \(f_1,g_1,h_1\). These marginals are Gaussian weights, so the induction hypothesis applies in dimension \(n\). Replace their integrals by those of the original weights to obtain

\[\left(\int f\right)^{1-t}\left(\int g\right)^t\le\int h.\]

The formal statement writes positive powers using exponentials and logarithms, including \(t=0\) and \(t=1\).




The induction splits \(\mathbb R^{d+1}\) into its first coordinate and \(d\) remaining coordinates. For fixed remaining coordinates, apply the proved one-dimensional inequality to the first-coordinate slices. Their integrals define marginal functions that satisfy the same interpolation inequality in \(d\) dimensions. Gaussian domination proves slice integrability, continuity of those marginal functions by dominated convergence, and Gaussian tails for them. The induction hypothesis applies, and Fubini identifies the iterated integrals with the original total integrals. Dimension zero is the one-point base case.




\textbf{Correspondence with the displayed declarations.}


\begin{itemize}

\item \texttt{prekopa\_\allowbreak{}leindler}: The finite-dimensional Prékopa--Leindler inequality. The induction uses the proved real-line transport theorem and integrates one coordinate at a time.
\end{itemize}


{\raggedright\noindent Lean declarations: \texttt{prekopa\_\allowbreak{}leindler}.\par}

\end{document}
